% =============================================================
%  Section: Results
% =============================================================

The results in this section are computed on the held-out COCO test split (5\,\%
of the consolidated dataset, $6\,165$ images, $42\,795$ ground-truth bounding
boxes after the removal of crowd annotations; cf.\ the discussion of the
one-object-per-cell drop in Section~\ref{sec:experiments:setup}). All
numbers correspond to the V6 checkpoint (\texttt{v6\_run/best\_v6.pth})
introduced in Section~\ref{sec:experiments:v6}. Inference uses per-class non-maximum
suppression with an IoU threshold of $0.5$ and is run at two prediction
probability thresholds, $0.1$ (permissive) and $0.3$ (stricter), in order to
characterise the precision/recall trade-off rather than commit to a single
operating point.

% -------------------------------------------------------------
\subsection{Headline metric}
\label{sec:results:headline}

Table~\ref{tab:headline} reports the mean Average Precision at IoU~$0.5$
(mAP@0.5) on the COCO test split for both probability thresholds. The
headline number for the project, used in the abstract and the conclusions, is
$\mathrm{mAP}@0.5 = 0.0415$ at probability threshold~$0.1$.

\begin{table}[!ht]
\centering
\caption{Final detection metrics of the V6 checkpoint on the COCO test split.
Averaged over the 80 COCO classes, all of which have at least one ground-truth
instance in the test split.}
\label{tab:headline}
\small
\begin{tabular}{lcc}
\toprule
\textbf{Probability threshold} & \textbf{mAP@0.5} & \textbf{Number of classes with AP $>$ 0} \\
\midrule
0.1 & \textbf{0.0415} & 60 / 80 \\
0.3 & 0.0381          & 58 / 80 \\
\bottomrule
\end{tabular}
\end{table}

\noindent We make no claim of competitiveness with current detectors; the
purpose of the figure is to bound the absolute performance of a paper-faithful
YOLOv1 implementation, re-targeted to COCO with a ResNet18 backbone, trained
for $\sim 60$ epochs in total on a single GPU. The remainder of this section
is dedicated to characterising \emph{where} the model spends and fails to
spend that detection budget.

% -------------------------------------------------------------
\subsection{Per-class breakdown}
\label{sec:results:perclass}

Of the 80 COCO classes evaluated at threshold~$0.1$, the AP distribution is
strongly skewed: the top decile (eight classes) accounts for the bulk of the
mAP value, while a long tail of twenty classes registers exactly zero. The
full ranked table for both thresholds is provided in
Appendix~\ref{app:per-class-ap}; we discuss the structure of the distribution
here.

\paragraph{Top performers.} Table~\ref{tab:top10} lists the ten classes with
the highest AP at probability threshold~$0.1$. The pattern is consistent:
classes with large, structurally distinctive silhouettes
(zebra stripes, giraffe necks, elongated train and airplane fuselages) and
classes that tend to be photographed as the main subject of the image
(cat, bear, elephant, bus) dominate the head of the distribution.

\begin{table}[!ht]
\centering
\caption{Top-10 classes by AP@0.5 (V6, threshold~$0.1$). The ``GT in test''
column is the number of ground-truth instances of the class in the test
split.}
\label{tab:top10}
\small
\begin{tabular}{lrrr}
\toprule
\textbf{Class} & \textbf{AP@0.5} & \textbf{GT in test} & \textbf{Class idx} \\
\midrule
zebra      & 0.3593 &  236 & 22 \\
train      & 0.3265 &  212 &  6 \\
giraffe    & 0.3107 &  225 & 23 \\
cat        & 0.2575 &  255 & 15 \\
airplane   & 0.1857 &  238 &  4 \\
pizza      & 0.1795 &  257 & 53 \\
bear       & 0.1698 &   67 & 21 \\
elephant   & 0.1634 &  223 & 20 \\
bus        & 0.1595 &  246 &  5 \\
toilet     & 0.1435 &  237 & 61 \\
\bottomrule
\end{tabular}
\end{table}

\paragraph{Classes with $\mathrm{AP}=0$.} At threshold~$0.1$, twenty COCO
classes register zero detections that pass mAP's IoU criterion. Inspecting
this set reveals four distinct failure modes, each of which is consistent with
the documented weaknesses of YOLOv1~\citep{redmon2016you} and is amplified by
our particular setting.

\begin{itemize}
    \item \emph{Objects smaller than one grid cell.} With $S=7$ and an input
        resolution of $448\times 448$ pixels, each grid cell covers a
        $64\times 64$ patch of the image. Toothbrushes, hair driers, scissors
        and remote controls occupy substantially less than this footprint in
        a typical photograph. YOLOv1's ``one object per cell'' constraint
        means the detection head has no way to localise multiple sub-cell
        objects, and the regression target for the single object it does
        attempt to localise is dominated by quantisation noise.
    \item \emph{Thin or elongated tableware in cluttered scenes.} Fork,
        knife, spoon and hot dog all appear most frequently in dining table
        and food-tray scenes where many of them are grouped within a single
        cell. The same constraint applies: only one of them can be
        ``responsible'' per cell.
    \item \emph{Rare classes.} Hair drier (5 GT in the test split), toaster
        (9), parking meter (59), scissors (78), toothbrush (80) and stop sign
        (76) have so few training instances that the model has had limited
        opportunity to learn a robust template for them. The exception in the
        top performers, bear (67 GT), shows that rarity alone is not
        sufficient to drive AP to zero; rarity \emph{combined with} a small
        or non-distinctive silhouette is.
    \item \emph{Sports objects in cluttered foregrounds.} Sports ball,
        frisbee, baseball bat, baseball glove and snowboard typically appear
        in scenes dominated by human figures with the object partially
        occluded or near the image periphery. The class loss for the
        responsible cell is therefore in direct competition with the much
        stronger gradient from the co-located ``person'' instance.
\end{itemize}

\paragraph{The frequency--AP decoupling.} A naive expectation is that AP
should grow monotonically with the number of ground-truth instances seen
during training. The data refutes this clearly. The ``person'' class is the
most common in COCO by a wide margin (8\,779 GT instances in the test split,
roughly 17\,\% of all annotations) and attains only $\mathrm{AP}=0.035$. The
``car'' class (1\,382 GT) attains $0.005$. ``Chair'' (1\,336 GT) attains
$0.002$. Conversely, zebra and giraffe, with under 240 instances each, are at
the top. The discriminating factor is not frequency but a combination of size,
distinctiveness and image context.

\paragraph{Effect of the probability threshold.} Moving from threshold~$0.1$
to $0.3$ drops the mAP by approximately $8\,\%$ (from $0.0415$ to $0.0381$).
Most of the absolute loss comes from classes whose model output is uncertain
even when correct: \emph{zebra} loses $2.2$ percentage points of AP and
\emph{giraffe} $1.2$ percentage points, while \emph{train} is essentially
unaffected ($-0.2$ points). The set of zero-AP classes grows by two
(\emph{mouse} and \emph{cell phone}, both flipping from $\sim\!10^{-4}$ to
$0$). The shift is consistent with a calibration in which true positives
at this stage of training are concentrated near the lower confidence values
rather than the higher ones.

% -------------------------------------------------------------
\subsection{Qualitative analysis}
\label{sec:results:qualitative}

We illustrate the model's behaviour with examples drawn from the test split.
The selection is partitioned into three subjective categories: \emph{good}
(every ground-truth object is detected with a class-correct, well-localised
box and few false positives), \emph{medium} (the primary subject is
recognised but localisation is imperfect or one of several objects is missed)
and \emph{bad} (the prediction is essentially uncorrelated with the
ground truth, or the only detected box is a false positive of a different
class).

\begin{figure}[!ht]
\centering
\begin{subfigure}[t]{0.32\textwidth}
    \includegraphics[width=\linewidth]{pred_015522_best_v6.png}
    \caption{Multiple distinct subjects.}
\end{subfigure}\hfill
\begin{subfigure}[t]{0.32\textwidth}
    \includegraphics[width=\linewidth]{pred_059397_best_v6.png}
    \caption{Single subject, tight box.}
\end{subfigure}\hfill
\begin{subfigure}[t]{0.32\textwidth}
    \includegraphics[width=\linewidth]{pred_089898_best_v6.png}
    \caption{High-AP class on a clean background.}
\end{subfigure}
\caption{Good detections (V6 checkpoint, probability threshold~$0.3$). The
common factor across the good set is a large, centrally placed subject from
one of the head classes of the distribution analysed in
Section~\ref{sec:results:perclass}.}
\label{fig:preds-good}
\end{figure}

\begin{figure}[!ht]
\centering
\begin{subfigure}[t]{0.32\textwidth}
    \includegraphics[width=\linewidth]{pred_001560_best_v6.png}
    \caption{Localisation drifts off the subject.}
\end{subfigure}\hfill
\begin{subfigure}[t]{0.32\textwidth}
    \includegraphics[width=\linewidth]{pred_016274_best_v6.png}
    \caption{Class correct, secondary object missed.}
\end{subfigure}\hfill
\begin{subfigure}[t]{0.32\textwidth}
    \includegraphics[width=\linewidth]{pred_079617_best_v6.png}
    \caption{Box too large for the object.}
\end{subfigure}
\caption{Medium-quality detections. The detector identifies the dominant
object in the scene but with imperfect localisation, or misses secondary
objects whose centres fall in cells already claimed by another instance.}
\label{fig:preds-medium}
\end{figure}

\begin{figure}[!ht]
\centering
\begin{subfigure}[t]{0.32\textwidth}
    \includegraphics[width=\linewidth]{pred_000001_best_v6.png}
    \caption{Class confusion.}
\end{subfigure}\hfill
\begin{subfigure}[t]{0.32\textwidth}
    \includegraphics[width=\linewidth]{pred_026429_best_v6.png}
    \caption{No detection above threshold.}
\end{subfigure}\hfill
\begin{subfigure}[t]{0.32\textwidth}
    \includegraphics[width=\linewidth]{pred_062225_best_v6.png}
    \caption{Cluttered scene with small objects.}
\end{subfigure}
\caption{Failure cases. Either the prediction is of a different class than
the ground truth (left), no prediction survives the probability filter
(middle), or the scene contains small objects below the effective resolution
of the $S=7$ grid (right).}
\label{fig:preds-bad}
\end{figure}

The qualitative split corroborates the per-class analysis: every example in
the good set features a single dominant subject from a head class, the medium
set fails on either secondary objects or imprecise localisation, and the bad
set is dominated by small-object and cluttered-scene cases. We did not curate
the bad examples to avoid the head classes; the bad failures simply occur
mostly outside the head.

% -------------------------------------------------------------
\subsection{Inference latency}
\label{sec:results:latency}

Although YOLOv1's original motivation was real-time inference, latency is not
the centrepiece of this work, because the training-time bug investigation in
Section~\ref{sec:bug-investigation} dominates the contribution. For
completeness, we report the forward-pass latency of the V6 checkpoint measured
by \texttt{scripts/realtime\_detect.py} in its \texttt{-{}-benchmark} mode,
which runs the model on synthetic batch-size-1 inputs after a warm-up and
records the per-call wall-clock time with explicit
\texttt{torch.cuda.synchronize} barriers before and after each forward pass.

\begin{table}[!ht]
\centering
\caption{Forward-only inference latency of the V6 model on an NVIDIA
GeForce RTX~5070~Ti~Laptop, batch size~1, $448\!\times\!448$ input. Statistics
over 50 timed iterations after a 20-iteration warm-up. Peak GPU memory
allocated: $\sim\!1.9$~GB.}
\label{tab:latency}
\small
\begin{tabular}{lrr}
\toprule
\textbf{Statistic} & \textbf{Latency (ms)} & \textbf{Throughput (FPS)} \\
\midrule
mean         & 5.32 & 188.1 \\
median       & 5.24 & 190.9 \\
$p_{95}$     & 6.27 & 159.5 \\
$p_{99}$     & 6.76 & 147.9 \\
\bottomrule
\end{tabular}
\end{table}

\noindent The model comfortably exceeds the real-time threshold typically
associated with the original paper, which reports 45~FPS for the base YOLO
model on a Titan~X. The gap is mostly explained by hardware
(an Ada-generation laptop GPU vs.\ a 2016-era Titan~X), not by the
architecture: the parameter count of our model (248~M) is in fact larger than
the paper's, but it is dominated by a single dense layer whose evaluation is
trivial on modern hardware. When the same script is run in its interactive
webcam mode---which adds frame capture, BGR$\to$RGB conversion, NMS, drawing
and a GUI window to the inference loop---throughput is correspondingly lower
and limited primarily by the webcam's frame-rate cap.

% -------------------------------------------------------------
\subsection{Comparison with the original paper}
\label{sec:results:vs-paper}

Table~\ref{tab:vs-paper} contrasts the configuration and headline result of
this work with the figures reported by \citet{redmon2016you}. The two are not
directly comparable as benchmarks---the datasets differ (COCO vs Pascal~VOC)
and so do the backbones and training budgets---but the comparison is useful
to bound the order-of-magnitude gap and to make explicit which contributing
factors are deliberate (dataset, backbone) and which are pragmatic (training
length, augmentation budget).

\begin{table}[!ht]
\centering
\caption{Side-by-side comparison with the original YOLOv1 paper.}
\label{tab:vs-paper}
\small
\begin{tabular}{lll}
\toprule
\textbf{Aspect} & \textbf{\citet{redmon2016you}} & \textbf{This work (V6)} \\
\midrule
Dataset                & VOC 2007 + 2012 (20 classes)   & COCO 2017 (80 classes) \\
Backbone               & Darknet24 (custom)             & ResNet18 (pretrained) \\
Backbone pretraining   & Manual on ImageNet             & torchvision IMAGENET1K\_V1 \\
Peak learning rate     & $10^{-2}$ (manual schedule)    & $10^{-4}$ \\
LR schedule            & Manual ($10^{-3}\!\to\!10^{-2}\!\to\!10^{-3}\!\to\!10^{-4}$) & Linear warm-up + cosine \\
Epochs                 & 135                             & 60 (V5 + V6 combined) \\
Batch size             & 64                              & 32 \\
Input resolution       & $448\times 448$                 & $448\times 448$ \\
Grid ($S$)             & 7                               & 7 \\
Boxes per cell ($B$)   & 2                               & 2 \\
Output channels per cell & $20 + 5\cdot 2 = 30$         & $80 + 5\cdot 2 = 90$ \\
Augmentation           & Crops, exposure, saturation, scaling & Resize only \\
$\lambda_{\text{coord}}$ & 5                            & 5 \\
$\lambda_{\text{noobj}}$ & 0.5                          & 0.5 \\
Hardware               & NVIDIA Titan~X (2016)          & NVIDIA H200~NVL (2024) \\
Wall-clock training    & $\sim$ 1 week                  & $\sim$ 2 hours \\
\midrule
\textbf{mAP@0.5}       & 63.4 (VOC 2007 test)            & 4.15 (COCO test, threshold~0.1) \\
\bottomrule
\end{tabular}
\end{table}

\noindent The gap of roughly $\times 15$ in headline mAP is dominated by four
factors that we discuss in detail in the Discussion: (i) COCO is a
substantially harder benchmark than VOC, both in terms of class count and in
terms of object size distribution, (ii) the original paper trains for more
than twice as long, (iii) ResNet18 is roughly $3\times$ smaller in parameter
count than Darknet24 in the backbone (although our overall model is larger
due to the FC head), and (iv) the augmentation pipeline used in the paper is
considerably richer than our resize-only transform.

We deliberately do not interpret the residual gap as evidence of an
implementation problem after V5/V6: the bug investigation in
Section~\ref{sec:bug-investigation} characterised the implementation bugs
exhaustively, the round-trip and synthetic verifications certified that the
loss and the coordinate decoder are now correct, and the remaining headroom
is the expected consequence of the four factors above. The Discussion section
sketches what would be required to close part of this gap within the same
architectural family.
